Após anos de dedicação ao curso, o aluno se depara com o desafio de confeccionar sua dissertação ou tese. Este documento visa orientar esse processo, esclarecendo as normas e boas práticas para a produção de trabalhos acadêmicos no \ac{ITA}.

Os modelos de trabalhos acadêmicos do \ac{ITA} seguem as normas da \ac{ABNT}, em especial a NBR 14724 (trabalhos acadêmicos) e a NBR 6023 (referências bibliográficas). Embora tais normas utilizem terminologia específica da área de documentação, este template busca simplificar sua aplicação através de formatação automatizada.

\section{Por que \LaTeX?}

O \LaTeX{} é um sistema de preparação de documentos amplamente utilizado na comunidade acadêmica e científica. Suas principais vantagens incluem:

\begin{enumerate}
    \item Qualidade tipográfica superior, especialmente para expressões matemáticas como $\int_0^\infty e^{-x^2} dx = \frac{\sqrt{\pi}}{2}$;
    \item Gerenciamento automático de referências cruzadas, numeração de figuras, tabelas e equações;
    \item Formatação consistente em todo o documento;
    \item Separação entre conteúdo e apresentação, permitindo que o autor foque no texto;
    \item Controle de versão facilitado por ser baseado em texto puro; e
    \item Ampla disponibilidade de pacotes para necessidades específicas.
\end{enumerate}

Este template utiliza o sistema Bib\LaTeX{} com \textit{backend} Biber para gerenciamento de referências bibliográficas, garantindo conformidade com as normas ABNT vigentes e oferecendo recursos avançados como \textit{backref} (indicação das páginas onde cada referência é citada).

\section{Tipos de Trabalhos Acadêmicos}

Este template atende aos seguintes tipos de trabalhos acadêmicos do \ac{ITA}:

\begin{enumerate}
    \item Trabalho de Graduação, para obtenção do diploma de Engenheiro;
    \item Dissertação de Mestrado, para obtenção do título de Mestre em Ciências;
    \item Dissertação de Mestrado Profissional, para obtenção do título de Mestre; e
    \item Tese de Doutorado, para obtenção do título de Doutor em Ciências.
\end{enumerate}

A configuração do tipo de documento é feita através das opções da classe, como \texttt{msc} para mestrado, \texttt{dsc} para doutorado, ou \texttt{tg} para trabalho de graduação.

\section{Estrutura do Documento}

A apresentação do texto deve obedecer à ordem dos elementos descritos na Tabela~\ref{tab:ordem-apresentacao}, conforme a ABNT NBR 14724.

\begin{table}[htbp]
\centering
\caption{Ordem de apresentação da Dissertação ou Tese.}
\label{tab:ordem-apresentacao}
\begin{tabular}{p{2.5cm}p{7.3cm}p{2.2cm}p{2.2cm}}
\toprule
Elemento & Estrutura (sequencial) & Versão preliminar & Versão final \\
\midrule
\multirow{15}{*}{Pré-textuais} & Capa & obrigatório & obrigatório \\
 & Folha de rosto & obrigatório & obrigatório \\
 & Verso da folha de rosto & não colocar & obrigatório \\
 & Folha da banca & obrigatório & obrigatório \\
 & Dedicatória & opcional & opcional \\
 & Agradecimentos & opcional & opcional \\
 & Epígrafe & opcional & opcional \\
 & Resumo & obrigatório & obrigatório \\
 & Abstract & obrigatório & obrigatório \\
 & Lista de Figuras & opcional & opcional \\
 & Lista de Tabelas & opcional & opcional \\
 & Lista de Abreviaturas e Siglas & opcional & opcional \\
 & Lista de Símbolos & opcional & opcional \\
 & Sumário & obrigatório & obrigatório \\
\midrule
\multirow{3}{*}{Textuais} & Introdução & obrigatório & obrigatório \\
 & Desenvolvimento & obrigatório & obrigatório \\
 & Conclusão & obrigatório & obrigatório \\
\midrule
\multirow{5}{*}{Pós-textuais} & Referências & obrigatório & obrigatório \\
 & Apêndice(s) & opcional & opcional \\
 & Anexo(s) & opcional & opcional \\
 & Glossário & opcional & opcional \\
 & Folha de Registro do Documento & não colocar & obrigatório \\
\bottomrule
\end{tabular}
\end{table}

Ao contrário da versão preliminar, na versão final devem ser incorporados os elementos ``Verso da folha de rosto'' e ``Folha de Registro do Documento''. A ``Folha de Registro do Documento'' não é um item da ABNT e sim um documento necessário para atender a NPA-ITA-011B:2023.

Tanto o Resumo quanto o Abstract são itens obrigatórios, independentemente se o texto seja escrito em português ou inglês.

\section{Estrutura da Introdução}

A introdução de uma dissertação ou tese deve conter os seguintes elementos:

\begin{enumerate}
    \item contextualização do tema e motivação para o estudo;
    \item revisão bibliográfica sucinta, situando o trabalho no estado da arte;
    \item definição clara dos objetivos (geral e específicos);
    \item no caso de doutorado, descrição das contribuições originais; e
    \item descrição da metodologia e organização dos capítulos.
\end{enumerate}

\section{Organização deste Documento}

Este documento está organizado da seguinte forma:

O Capítulo 2 apresenta orientações sobre revisão bibliográfica e demonstra o uso de citações e figuras.

O Capítulo 3 exemplifica um capítulo de Metodologia, comum em dissertações e teses.

O Capítulo 4 aborda aspectos de formatação, incluindo margens, fontes, espaçamento e paginação conforme as normas ABNT.

O Capítulo 5 trata da elaboração de citações e referências bibliográficas, demonstrando os diversos tipos de fontes e suas formatações.

O Capítulo 6 orienta sobre a confecção da versão final em mídia digital para depósito na biblioteca.

O Capítulo 7 explica o preenchimento da Folha de Registro do Documento.

O Capítulo 8 apresenta as considerações finais.
